\section{Method}
\label{sec:method}

The methodology's central idea is a strict division of labor between
three actors, mediated by a version-controlled repository that is the
single source of truth for the design
(Fig.~\ref{fig:workflow}):

\begin{enumerate}[nosep,label=(\roman*)]
\item \textbf{Deterministic physics code} is the \emph{only} source of
  engineering numbers. Every sizing relation, mass model, and margin
  check is transparent, configured Python under regression test. No
  quantity used in a design decision is ever generated by the language
  model itself.
\item \textbf{The AI agent} --- an LLM-based agentic coding tool
  \citep{anthropic2025claudecode} with full read, write, and execute
  access to the repository --- writes, refactors, and re-converges
  that code,
  propagates design changes across every affected discipline, drafts
  decision records and documentation, and executes the analysis
  pipeline. It is an engineering \emph{collaborator}, not an oracle:
  its contributions are ordinary commits subject to the same review
  and testing gates as human-written ones.
\item \textbf{Human engineers} set requirements, issue design
  directives in natural language, review every change (via pull
  requests), perform independent expert review of the evolving design,
  and accept or reject each recorded decision.
\end{enumerate}

Sections~\ref{sec:pipeline}--\ref{sec:iteration} describe the four
pillars that make this collaboration rigorous: a one-way analysis
pipeline, a machine-readable governance layer, a verification harness
that the AI cannot silently bypass, and an interaction protocol with
full provenance.

\begin{figure}[htbp]
\centering
\resizebox{\textwidth}{!}{%
\begin{tikzpicture}[
  font=\small,
  node distance=5mm and 9mm,
  box/.style={draw, rounded corners=2pt, align=center, minimum height=8mm,
              inner sep=2.5mm, fill=gray!8},
  actor/.style={draw, align=center, minimum height=8mm, inner sep=2.5mm,
                fill=blue!8, rounded corners=2pt},
  gate/.style={draw, align=center, minimum height=8mm, inner sep=2.5mm,
               fill=orange!12, rounded corners=2pt},
  arr/.style={-{Stealth[length=2.2mm]}, thick}
]
% pipeline row
\node[box] (config) {\code{config/*.yaml}\\design inputs};
\node[box, right=of config] (src) {\code{src/}\\physics library};
\node[box, right=of src] (nb) {\code{notebooks/}\\16-stage pipeline};
\node[box, right=of nb] (out) {\code{out/}\\typed hand-offs};
\node[box, right=of out] (cfd) {CAD / CFD /\\flight-stack\\consumers};
\draw[arr] (config) -- (src);
\draw[arr] (src) -- (nb);
\draw[arr] (nb) -- (out);
\draw[arr] (out) -- (cfd);
% actors row
\node[actor, above=11mm of src] (human) {Human engineers\\directives, review,\\expert findings};
\node[actor, above=11mm of out] (ai) {AI agent\\implements, re-converges,\\documents};
\draw[arr, <->] (human) -- node[above]{natural language} (ai);
\draw[arr] (ai.south) -- ++(0,-4mm) -| (nb.north);
\draw[arr] (human.south) -- ++(0,-4mm) -| (config.north);
% gates row
\node[gate, below=11mm of src] (gov) {Governance\\constitution (\code{CLAUDE.md}),\\ADRs, changelog};
\node[gate, below=11mm of out] (ver) {Verification\\design-point pins, CI\\full-pipeline re-execution};
\draw[arr, <->] (gov) -- (ver);
\draw[arr] (gov.north) -- ++(0,4mm) -| (src.south);
\draw[arr] (ver.north) -- ++(0,4mm) -| (out.south);
\end{tikzpicture}%
}
\caption{The governed human--AI--repository loop. Engineering numbers
flow only left-to-right through the deterministic pipeline (top row);
the AI agent and human engineers modify the repository under the
governance and verification gates (bottom row).}
\label{fig:workflow}
\end{figure}

\subsection{The design repository as a one-way analysis pipeline}
\label{sec:pipeline}

All design knowledge lives in one repository with a strictly one-way
dataflow:
\begin{equation*}
\code{config/} \;\rightarrow\; \code{src/} \;\rightarrow\;
\code{notebooks/} \;\rightarrow\; \code{out/} \;\rightarrow\;
\text{CAD/CFD/flight-stack consumers.}
\end{equation*}
Every physical parameter --- mission times, aerodynamic limits,
battery properties, structural factors --- lives in commented YAML
configuration files with a written rationale; the Python physics
library contains \emph{no} magic numbers. Jupyter notebooks are thin
orchestration only (parameters, module calls, interpretation); all
physics, plotting, and reporting code lives in the importable library,
so the entire design is reproducible headlessly and testable by unit
tests. Each pipeline stage writes typed YAML/JSON hand-off files that
downstream stages read, and no stage ever writes back upstream. In the
case study this pipeline comprises sixteen notebook stages
(Table~\ref{tab:pipeline}), executed in dependency order.

This architecture is not incidental to the AI methodology --- it is
what makes an LLM agent safe to employ. Because the dataflow is
one-way and fully regenerable from \code{config/}, the agent can be
asked to re-converge the whole design after any change, and the effect
of the change is completely observable in the diff of the generated
hand-offs. Because notebooks are thin, the agent's edits concentrate
in ordinary, reviewable, unit-testable Python modules rather than in
opaque notebook state.

\begin{table}[htbp]
\centering
\caption{The sixteen-stage analysis pipeline of the case study. Each
stage writes typed hand-offs to \code{out/} that later stages consume;
stages 13--15 re-solve stages 4--6 with the frozen commercial hardware
without touching the conceptual outputs (no back-edge).}
\label{tab:pipeline}
\small
\begin{tabularx}{\textwidth}{@{}rlX@{}}
\toprule
\# & \textbf{Stage} & \textbf{Function} \\
\midrule
1 & Conceptual design & Mission sizing and iterative mass closure (MTOW, wing, power) \\
2 & Wing design & Airfoil selection and wing geometry \\
3 & Control vane design & Jet-vane sizing from hover thrust \\
4 & Aileron design & Cruise-phase roll authority (backup to vanes) \\
5 & Vibration isolation & FC/IMU and payload soft mounts vs.\ 1/rev rotor forcing \\
6 & Fuselage design & Packaging, center of gravity, drag trade \\
7 & Thermal design & ESC cold plate, vented battery bay, pack mission transient \\
8 & Vehicle solid model & Parametric CAD (STEP/STL export) \\
9 & Aerodynamic hand-off & Versioned JSON geometry contract for external VLM/BEMT analysis \\
10 & Mass properties & Inertia tensor, bill of materials, CG trim \\
11 & Wiring diagram & Generated electrical block diagram from the sizing result \\
12 & COTS selection & Deterministic component freeze against derived requirements \\
13--15 & As-selected re-solve & Stages 4--6 re-solved with frozen hardware masses/envelopes \\
16 & Design summary & Pure reader: design point, margins, standing findings \\
\bottomrule
\end{tabularx}
\end{table}

\subsection{Governance: a machine-readable project constitution}
\label{sec:governance}

The AI agent is steered by a \emph{project constitution} --- a
repository file (\code{CLAUDE.md}) that the agent loads at the start
of every working session and that encodes the engineering rules of the
program in plain language. In the case study it fixes, among others:

\begin{itemize}[nosep]
\item \textbf{No magic numbers}: every physical parameter lives in
  configuration with a commented rationale; numerical tolerances and
  named guard limits are the only exceptions.
\item \textbf{All physics in the library}: notebooks only set
  parameters, call modules, and interpret results.
\item \textbf{Conventions}: a single body-fixed axis convention
  (front--right--down) and unit discipline (SI in hand-offs and module
  interfaces; millimetres only inside CAD exports) applied uniformly
  across sizing, CAD, CFD, and documentation.
\item \textbf{Analysis-cost policy}: continuous integration validates
  CFD case \emph{setup} with a coarse smoke run only; converged CFD is
  never run in CI.
\item \textbf{Documentation duties}: every architecturally significant
  decision gets an architecture decision record (ADR); every tagged
  version gets a changelog entry.
\end{itemize}

Because the constitution is versioned alongside the code, the rules
evolve by the same review process as the design itself, and a new
human collaborator or a fresh AI session inherits identical
constraints. Design decisions are captured at decision time as ADRs
(status/context/decision/consequences); superseded decisions are never
rewritten but amended by new records, preserving the full reasoning
trail. In the case study the AI drafted each ADR as part of
implementing the corresponding change, and the human engineers
accepted or revised it in review --- eighteen ADRs over the program,
covering choices from ``all physics in \code{src/}'' to the propeller
blade section.

\subsection{Verification: making silent design drift impossible}
\label{sec:verification}

Three mechanisms ensure that neither the AI nor a human can move the
design unintentionally:

\paragraph{Design-point regression pins.} The converged design point
--- masses, geometry, control-surface torques, CFD reference values
--- is pinned in the test suite (147 test functions in the case
study). An \emph{intentional} design change must update the pins in
the same commit, which makes the diff itself the documentation of the
new design point; an \emph{unexpected} pin failure means the design
drifted and blocks the merge. This converts the classic LLM risk of
plausible-but-wrong output into a loud, reviewable event.

\paragraph{Full-pipeline continuous integration.} Every pull request
re-executes all sixteen notebook stages from scratch, runs the
regression and geometry tests (e.g., watertightness and unit checks on
exported CAD), and archives the regenerated hand-offs. On the main
branch, a coarse CFD smoke run additionally validates the OpenFOAM
case setup against the current geometry. Pushing a version tag
freezes a complete design snapshot (YAML hand-offs, STEP/STL geometry,
rendered notebooks) as an immutable release.

\paragraph{Cross-consumer consistency checks.} Values that must agree
across tool boundaries --- e.g., the CFD case's reference length,
area, and center of rotation versus the fuselage hand-off --- are
cross-checked by tests, so a design move that forgets a consumer fails
loudly. During the case study these checks caught, for example,
propulsion- and vane-CFD geometry that had silently drifted from the
converged design point, which was then re-sourced from the hand-offs
at run time.

A further discipline governs \emph{negative} results: margins that
come out unfavorable are recorded as pinned ``standing findings''
(reported by the final pipeline stage on every run) rather than tuned
away. Section~\ref{sec:results} lists the case study's standing
findings; methodologically, the point is that the AI is instructed to
surface unfavorable physics honestly, and the verification layer
prevents quiet re-tuning.

\subsection{Interaction protocol and provenance}
\label{sec:protocol}

Work proceeds in conversational sessions between an engineer and the
AI agent. The engineer's inputs are \emph{mission-level directives},
not code: representative examples from the case study include
``re-select the airfoil now that the stall-limited sizing is
honest,'' ``adopt a Clark~Y section for the propeller blade and
determine the rotation sense,'' ``fold the battery pack's internal
resistance into the closure consistently,'' and pasted findings from
an external design review. For each directive the agent:

\begin{enumerate}[nosep]
\item locates every discipline the change touches (sizing, geometry,
  control, thermal, CAD, CFD references, documentation);
\item implements the change in the physics library and configuration,
  with rationale comments;
\item re-executes the pipeline and re-converges the design point;
\item updates the regression pins, downstream consumer references, the
  ADR, and the changelog \emph{in the same commit}; and
\item opens a pull request that the human reviews and merges.
\end{enumerate}

Two provenance mechanisms make the AI's role auditable. First, every
AI-implemented commit carries a machine-readable co-authorship trailer
identifying the model, so the exact human/AI contribution split is
recoverable from the git history (Section~\ref{sec:process-results}).
Second, refactoring tasks are verified by \emph{byte-identical output
checks}: when the agent moved all notebook physics into the library
(ADR-0013 in the case study) and later split a pipeline stage
(ADR-0018), the generated hand-offs were verified byte-for-byte
unchanged before and after --- a strong guarantee that a structural
change carried no hidden design change.

The agent also maintains a persistent, human-readable memory of
project decisions and preferences across sessions (e.g., procurement
constraints such as preferring domestic retailers, or the standing
instruction to delete stale generated artifacts), complementing the
constitution with accumulated project context.

\subsection{Evidence-based iteration and external review}
\label{sec:iteration}

The design evolves through short, complete iterations: every
requirement change, model-fidelity improvement, or external finding
triggers a full re-convergence and a tagged snapshot, so the program
history is a sequence of internally consistent design points rather
than a single drifting draft. Three evidence channels feed this loop:

\begin{itemize}[nosep]
\item \textbf{Model refinement}: replacing an assumption with physics
  (e.g., an explicit structural member model replacing an area-density
  estimate; mission-transient battery heating replacing a steady-state
  vent check) and letting the closure re-converge honestly, even when
  the design gets heavier.
\item \textbf{Reality intrusion via procurement}: a deterministic
  commercial-component selection stage derives hard requirements from
  the converged design point and freezes the lightest feasible
  candidates; where no purchasable part matches an assumption (as
  happened with the case study's original fan diameter), the
  assumption --- not the report --- is amended, by ADR, and the design
  re-converged. A subsequent as-selected re-solve stage quantifies the
  gap between the conceptual closure and the frozen hardware without
  contaminating the conceptual baseline.
\item \textbf{Independent expert review}: domain experts outside the
  daily loop review the evolving design; their findings enter as
  directives and are traceable to the resulting ADRs. In the case
  study, external review contributed the semi-monocoque fuselage
  architecture, exposed a stall-speed inconsistency the automated
  checks could not see (a requirement silently unmet because a
  placeholder coefficient survived in one card), and prompted the
  battery internal-resistance accounting fix --- each of which moved
  the design point (Section~\ref{sec:results}).
\end{itemize}

Finally, the conceptual repository hands off to higher-fidelity
consumers through versioned, schema-validated contracts: a JSON
geometry hand-off (planform, body of revolution, control surfaces,
blade stations, moment reference point) feeds an external
vortex-lattice/blade-element analysis loop, while exported CAD and the
OpenFOAM cases support CFD on the same geometry. Schema evolution is
itself governed --- each field addition in the case study was a
reviewed, versioned, backward-compatible schema release, several of
them fixing real consumer-side integration faults (e.g., a missing
wing--body placement anchor that had put wing control points inside
the fuselage in the consumer's solver).
